\paragraph{Corollary 5.5.}
Under the same hypothesis of Lemma 5.4, it holds that
\[rKt^D(x)\leq m+O(\log^3 n).\]
\begin{proof}
By an averaging argument, there exists a string $z\in\{0,1\}^d$ such that
\[\Pr_w[x\in R^D(A(x,z),z,w)]\geq\frac1{2m^2}.\]
Let $\alpha:=A(x,z)\in\{0,1\}^m$. The string $x$ can be described by the following pseudo-deterministic algorithm $M$: $M$ takes the advice string $\alpha\in\{0,1\}^m$ and $z\in\{0,1\}^d$ as input, picks $w\sim\{0,1\}^r$ randomly, and outputs a random element of $R^D(\alpha,z,w)$. Since the list size is at most $L=\operatorname{poly}(m)$, the probability that the output of $M$ is equal to $x$ is at least $\lambda:=\frac1{2m^2}\cdot\frac1L$. Therefore, we obtain $rKt_\lambda(x)\leq m+d+O(\log n)$. The result follows from Lemma 5.2.
\end{proof}
\paragraph{Lemma 5.6.}
If $\mathsf{MrKtP}\in\mathsf{BPTIME}(2^{\epsilon n})$, then for all sufficiently large $n\in\mathbb N$ and any strings $x\in\{0,1\}^n$ and $y\in\{0,1\}^n$, it holds that
\[rKt(x\mid y)+rKt(y)\leq rKt(x,y)+O(\epsilon n).\]
\begin{proof}
We may assume without loss of generality that $rKt(y)\gg\log^3 n$. Let $m$ and $m'$ be parameters chosen later. Let $G:\{0,1\}^n\times\{0,1\}^d\to\{0,1\}^m$ and $G':\{0,1\}^n\times\{0,1\}^{d'}\to\{0,1\}^{m'}$ be the pseudorandom generator construction of Lemma 5.4 with parameters $m$ and $m'$, respectively.
Fix $z\in\{0,1\}^d$ and $z'\in\{0,1\}^{d'}$. Since the strings $(G(x;z),G'(y;z'))$ can be described by $z,z'$, and a program that computes $(x,y)$, we obtain
\begin{multline*}
rKt(G(x;z),G'(y;z'))\leq rKt(x,y)+|z|+|z'|+O(\log n)\\
=rKt(x,y)+O(\log^3 n).\tag{13}
\end{multline*}
Pick $w\sim\{0,1\}^m$ and $z'\sim\{0,1\}^{d'}$ randomly. We claim that $rKt(w,G'(y;z'))\geq|w|+|w'|-O(\log n)$ with high probability. Toward a contradiction, assume that
\[\Pr_{w,z'}[rKt(w,G'(y;z'))<|w|+|w'|-O(\log n)]\geq\frac12.\]
By Fact 3.2, we have
\[\Pr_{w,w'}[rKt(w,w')\geq|w|+|w'|-O(\log n)]\geq1-o(1).\]
We define an oracle $D$ so that $D(w';w):=1$ if and only if $rKt(w,w')<|w|+|w'|-O(\log n)$. Then, the two inequalities above show that $D$ distinguishes the output distribution of $G'(y;-)$ from the uniform distribution. By Corollary 5.5, we obtain
\[rKt^D(y)\leq m'+O(\log^3 n).\]
Using the assumption that $\mathsf{MrKtP}\in\mathsf{BPTIME}(2^{\epsilon n})$, the oracle $D$ can be computed by a randomized algorithm running in time $2^{O(\epsilon m')}$. It follows that
\[rKt(y)\leq m'+O(\epsilon m')=(1+O(\epsilon))\cdot m'.\]
We choose $m'$ so that this inequality does not hold; that is,
\[m':=(1-O(\epsilon))\cdot rKt(y).\]
Then, we obtain a contradiction and conclude that
\[\Pr_{w,z'}[rKt(w,G'(y;z'))\geq m+m'-O(\log n)]\geq\frac12.\tag{14}\]
Define $D_y$ to be an oracle such that $D_y(w;z'):=1$ if and only if $rKt(w,G'(y;z'))\leq rKt(x,y)+O(\log^3n)$. By Equations (14) and (13), $D_y$ distinguishes the output distribution of $G(x;-)$ from the uniform distribution if
\[m+m'-O(\log n)\geq rKt(x,y)+O(\log^3 n).\]
We define $m:=rKt(x,y)-m'+O(\log^3n)$ so that this inequality holds. By Corollary 5.5, we obtain
\[rKt^{D_y}(x)\leq m+O(\log^3n).\]
Since the oracle $D_y$ can be computed in time $2^{O(\epsilon m)}$ given $y$ as hard-wired input, we conclude that
\begin{align*}
rKt(x\mid y)&\leq m+O(\epsilon m).\\
&\leq rKt(x,y)-m'+O(\epsilon n).\\
&\leq rKt(x,y)-(1-O(\epsilon))\cdot rKt(y)+O(\epsilon n).\\
&\leq rKt(x,y)-rKt(y)+O(\epsilon n).
\end{align*}
\end{proof}
\begin{proof}[Proof of Theorem 5.3]
The outline of the proof is as follows: First, we find a string $y\in\{0,1\}^n$ such that $rKt(y)\geq n$ by an exhaustive search over all strings $y\in\{0,1\}^n$. Next, we find a string $x\in\{0,1\}^n$ such that $rKt(x,y)\geq2n-O(\epsilon n)$ by an exhaustive search over all strings $x\in\{0,1\}^n$. The existence of such a string $x$ is guaranteed by Lemma 5.6. Finally, we observe that the pair $(x,y)$ can be pseudo-deterministically computed in time $2^{n+O(\epsilon n)}$, which contradicts the lower bound on $rKt(x,y)$. Details follow.
Using the assumption that $\mathsf{MrKtP}\in\mathsf{BPTIME}(2^{\epsilon n})$, let $M$ be the randomized algorithm that computes $rKt(y)$ on input $y$ of length $n$ in time $2^{O(\epsilon n)}$ with probability at least $1-2^{-2n}$. Fix $n\in\mathbb N$. Let $y_n$ be the lexicographically first string $y\in\{0,1\}^n$ such that $rKt(y)\geq n$. Note that the existence of such a string $y_n$ is guaranteed by Fact 3.2.
We claim that there exists a pseudo-deterministic algorithm $M_1$ that computes $y_n$ on input $n\in\mathbb N$ in time $2^{n+O(\epsilon n)}$. For all strings $y\in\{0,1\}^n$ in lexicographical order, the algorithm $M_1$ tests whether $rKt(y)\geq n$ using $M$, and outputs the first string $y$ that passes this test. Since the error probability of $M$ is at most $2^{-2n}$, by a union bound, the algorithm $M_1$ computes $y_n$ pseudo-deterministically with probability $1-2^{-n}$. The running time of $M_1$ is at most $2^{n+O(\epsilon n)}$.
Next, let $x_n$ be the lexicographically first string $x\in\{0,1\}^n$ such that
\[rKt(x,y_n)\geq2n-c\cdot\epsilon n,\tag{15}\]
where $c$ is a large constant to be chosen independently of $\epsilon$. We prove that $x_n$ is well defined, i.e., there exists a string $x\in\{0,1\}^n$ such that $rKt(x,y_n)\geq2n-c\cdot\epsilon n$. By Fact 3.2, there exists a string $x\in\{0,1\}^n$ such that $rKt(x\mid y)\geq n$. By Lemma 5.6,
\[rKt(x,y)+O(\epsilon n)\geq rKt(x\mid y)+rKt(y)\geq2n.\]
Choosing a large constant $c$, we obtain $rKt(x,y)\geq2n-c\cdot\epsilon n$, as desired.
We claim that there exists a pseudo-deterministic algorithm $M_2$ that computes $x_n$ given $y_n$ as input in time $2^{n+O(\epsilon n)}$. For all strings $x\in\{0,1\}^n$ in lexicographical order, the algorithm $M_2$ tests whether $rKt(x,y_n)\geq2n-c\cdot\epsilon n$ using $M$, and outputs the first string $x$ that passes this test. Using a union bound, the error probability of $M_2$ is bounded by $2^{-n}$. The running time of $M_2$ is at most $2^{n+O(\epsilon n)}$.
Finally, by combining the algorithms $M_1$ and $M_2$, we obtain an algorithm that computes $(x_n,y_n)$ pseudo-deterministically in time $2^{n+O(\epsilon n)}$. By the definition of $rKt$, this implies that
\[rKt(x_n,y_n)\leq n+O(\epsilon n).\]
We also have $rKt(x_n,y_n)\geq2n-O(\epsilon n)$ by Equation (15). Therefore, we obtain $2n-O(\epsilon n)\leq n+O(\epsilon n)$, which is a contradiction for a sufficiently small constant $\epsilon>0$.
\end{proof}
We observe that $\mathsf{MrKtP}$ is in the exponential-time variant of $\mathsf{PP}$.
\paragraph{Proposition 5.7.}
$\mathsf{MrKtP}\in\mathsf{PEXP}$.
\begin{proof}
Since $\mathsf{PP}$ is closed under truth-table reductions [26], it suffices to show that $\mathsf{MrKtP}$ is reducible to $\mathsf{PP}$ via an exponential-time truth-table reduction. Let $(x,1^s)$ be an input such that $x\in\{0,1\}^n$ and $s\in\mathbb N$. For each string $d\in\{0,1\}^*$ of length at most $s$, we ask the $\mathsf{PP}$ oracle whether $U(d,r)$ outputs $x$ in time $t$ for at least a $2/3$-fraction of $r\in\{0,1\}^t$, where $t:=2^{s-|d|}$. We accept the input $(x,1^s)$ if and only if a positive answer is returned from the oracle.
\end{proof}
We expect that Theorem 5.3 can be extended to a lower bound against $\mathsf{PP}$-type algorithms. However, the original motivation of Oliveira [62] is to study a gap version of $\mathsf{MrKtP}$, for which the upper bound of $\mathsf{pr\text{-}BPE}$ can be proved. Specifically, let $\mathsf{GapMrKtP}$ denote the promise problem $(\Pi_{\mathrm{Yes}},\Pi_{\mathrm{No}})$ such that $\Pi_{\mathrm{Yes}}=\{(x,s):rKt(x)\leq s\}$ and $\Pi_{\mathrm{No}}=\{(x,s):rKt(x)>s+c\log|x|\}$, where $c$ is a sufficiently large constant. One can observe $\mathsf{GapMrKtP}\in\mathsf{pr\text{-}BPE}$ [62]. Theorem 5.3 can be extended to this promise problem.
Here, abusing a notation, for a class $C$ of languages, we say that a promise problem $\Pi\in C$ if there exists a language $L\in C$ such that $\Pi_{\mathrm{Yes}}\subseteq L\subseteq\{0,1\}^*\setminus\Pi_{\mathrm{No}}$ for $(\Pi_{\mathrm{Yes}},\Pi_{\mathrm{No}})=\Pi$.
The lower bound of Theorem 5.8 does work against a bounded probabilistic algorithm that satisfies the promise of $\mathsf{BPP}$-type algorithms over all the inputs. In contrast, the upper bound $\mathsf{pr\text{-}BPE}=\mathsf{pr\text{-}BPTIME}(2^{O(n)})$ holds only for an algorithm that may not satisfy the promise of $\mathsf{BPP}$-type algorithms for some inputs. This leaves a qualitative gap between the lower bound and the upper bound on $\mathsf{GapMrKtP}$; closing this gap is an interesting open question. We note, however, that there is no gap under the plausible assumption that $\mathsf E\not\subseteq\mathsf{i.o.SIZE}(2^{\epsilon n})$ for some constant $\epsilon>0$, in which case $\mathsf{pr\text{-}BPP}=\mathsf P=\mathsf{BPP}$ [46].
\begin{proof}[Proof of Theorem 5.8]
The proof is essentially the same with Theorem 5.3. We only explain how to modify the proof of Theorem 5.3.
Toward a contradiction, we assume the existence of a randomized algorithm $M$ that computes some language $L$ infinitely often, where $L$ is a language consistent with $\mathsf{GapMrKtP}$. For all large $n\in\mathbb N$ and every $x\in\{0,1\}^n$ and $s\in[n]$, we assume that the binary encoding of $(x,s)$ is exactly equal to $m(n)$, where $m(n)=\Theta(n)$ is some function. This ensures that for infinitely many $n\in\mathbb N$, the algorithm $M$ computes $L$ correctly on input $(x,s)$ for every $x\in\{0,1\}^n$ and every $s\in[n]$.
First, we explain how to extend the lower bound of Theorem 5.3 to the infinitely-often version. Let $n$ be the length of $x_n\in\{0,1\}^n$ and $y_n\in\{0,1\}^n$ constructed in the proof of Theorem 5.3. By inspection, in the proof of Theorem 5.3, we use the algorithm $M$ on inputs of length at most $n'$, where $n'=O(n)$. Any input $x$ of length less than $n'$ can be padded to an input of length $n'$ by mapping $x$ to $x':=x1^{n'-|x|}$. This ensures that the length of any input to $M$ that we use in the proof is exactly equal to $n'$. Moreover, the correctness of the proof remains unchanged because $rKt(x)=rKt(x')\pm O(\log n)$.
Second, we explain how to deal with the promise problem. In the proof of Theorem 5.3, we exhaustively search the lexicographically first string $y_n\in\{0,1\}^n$ such that $rKt(y_n)\geq n$. Instead, we search the lexicographically first string $y_n\in\{0,1\}^n$ such that $(y_n,n-2c\cdot\log n)\notin L$. Since $L\subseteq\{0,1\}^*\setminus\Pi_{\mathrm{No}}$, this ensures that $rKt(y_n)\geq n-c\log n$. Similarly, we search the lexicographically first string $x_n\in\{0,1\}^n$ such that $((x_n,y_n),2n-O(\epsilon n))\notin L$. This ensures that $rKt(x_n,y_n)\geq2n-O(\epsilon n)$. Since the language $L$ can be decided by the randomized algorithm $M$, we can compute the pair $(x_n,y_n)$ pseudo-deterministically in time $2^{n+O(\epsilon n)}$. This is a contradiction.
\end{proof}
